\chapter{Datasets}
\label{ch:datasets}
% Opening paragraph: four datasets, two synthetic (ours), two from the
% literature. Only LOFAR has real telescope data with human labels -- say up
% front that this is the one that matters most.

\section{Synthetic dataset}
\label{sec:synthetic}
% Write about your generator. Sources: PART 8, PART 9, "Verified facts"
% section of the context doc, Synthetic Dataset*/dataset_statistics.txt.

\subsection{Generation}
Real radio data with reliable pixel-level RFI labels is scarce, so the first
part of this work used a synthetic dataset that I generated myself. The
generator is written in Python and uses NumPy's random number generator with a
fixed seed (42), so the whole dataset can be regenerated bit for bit.

Each image is a spectrogram: one axis is frequency and the other is time, and
each pixel holds the power received in that frequency channel at that moment.
The background is not plain random noise. It is built from a sky pedestal that
varies across frequency, a smooth astronomical signal, a slow drift in time,
and Gaussian noise on top. RFI is then injected into this background. Six
types of RFI are modelled, chosen to cover the shapes commonly seen in real
observations (\cref{tab:rfi_types}). Each image receives a random selection of
types at random strengths, and images are spread over four density levels:
clean (8.2\%), light (30.0\%), moderate (36.2\%) and heavy (25.6\%).

\begin{table}[H]
  \centering
  \caption{The six RFI types in the synthetic dataset, with the number of the
    1000 images (1024$\times$265 version) that contain each type.}
  \label{tab:rfi_types}
  \begin{tabular}{llr}
    \toprule
    Type & Appearance in the spectrogram & Images \\
    \midrule
    Narrowband      & thin line at one frequency, lasting for some time  & 889 \\
    Scattered       & small isolated hits scattered over the image        & 899 \\
    Blob            & patchy, cloud-like region                           & 673 \\
    Persistent band & group of channels affected for the whole observation & 588 \\
    Broadband block & rectangle covering many channels for a while        & 488 \\
    Wideband burst  & short flash across many channels at once            & 310 \\
    \bottomrule
  \end{tabular}
\end{table}

The dataset contains 1000 images, split into 700 for training, 150 for
validation and 150 for testing. I checked that no image appears in more than
one split by comparing SHA-256 hashes of the raw pixels; there are no shared
hashes.

\subsection{Deciding which pixels are RFI}
\label{sec:synthetic_label}
Because the RFI is injected by the generator, its exact strength at every
pixel is known. A pixel is labelled as RFI when the injected RFI power there
is greater than 0.5 times the standard deviation of the local noise. Pixels
where RFI was added but is weaker than this are labelled clean, because at
that level there is effectively nothing left in the data for a detector to
find.

This rule gives labels that are noise-free and unambiguous: every labelled
pixel is, by construction, measurably above the noise. In the 276$\times$600
version, 90.9\% of the RFI pixels lie at or above one noise standard
deviation. This makes the task considerably easier than labelling real data,
a point that becomes important in \cref{ch:drivers}.

The rule was refined between the two versions of the dataset. In the
276$\times$600 version a pixel was labelled RFI if it lay inside the geometric
outline of an injected RFI shape \emph{or} exceeded the 0.5$\sigma$ threshold.
Because this was a union, it could mark pixels that carried almost no RFI
power (0.06\% of labelled pixels in the training split, and up to 4.49\% in
the worst single image). In the 1024$\times$265 version only the threshold is
used, and it is applied \emph{after} the instrument response
(\cref{sec:bandpass}) and against the total local noise of each channel, so
RFI that the instrument has genuinely buried below the noise is not
labelled.

\subsection{Two versions of the dataset}
\label{sec:synthetic_versions}
Two versions of the dataset were used (\cref{tab:synth_versions}). My first
images were 1024$\times$1024 pixels. At that size only one image fitted in a
training batch on the 6\,GB GPU, which made training unstable. Akeret et
al.\footcite{akeret2017} trained \tfunet{} on images of 276$\times$600 pixels, so
I regenerated the dataset at exactly that size to reproduce their setup as
closely as possible. When the image size changes, the RFI shapes have to be
rescaled as well: a blob 40 channels wide covers 3.9\% of a 1024-channel band
but 14.5\% of a 276-channel band. Seven size-dependent parameters of the
generator were therefore scaled in proportion to the image dimensions.

The second version, 1024$\times$265, was made on my supervisor's advice for
two reasons. First, RFI is often very faint and occupies only a few frequency
channels, so a detector benefits from finer frequency resolution. The new
version has 1024 frequency channels instead of 276, which is 3.7 times finer
in frequency; in exchange it has 265 time samples instead of 600. Second, real
receivers do not respond equally to all frequencies, and this instrument
bandpass was missing from the first version. The second version adds it, as
described next.

\begin{table}[H]
  \centering
  \caption{The two versions of the synthetic dataset. Both use seed 42 and a
    700/150/150 split.}
  \label{tab:synth_versions}
  \small
  \begin{tabular}{lccc}
    \toprule
    Version & Size (freq $\times$ time) & Instrument bandpass & RFI pixels \\
    \midrule
    v3 & 276$\times$600  & no  & 14.67\% mean (test split 15.08\%) \\
    v4 & 1024$\times$265 & yes & 14.19\% mean \\
    \bottomrule
  \end{tabular}
\end{table}

\subsection{The instrument bandpass (v4)}
\label{sec:bandpass}
A radio receiver does not amplify every frequency by the same amount. Its gain
falls away at the edges of the observing band, and across the band it has
small ripples caused by the filters and by reflections between parts of the
telescope (standing waves). This frequency-dependent gain is called the
bandpass. Real spectrograms always carry it, so I added it to the second
version of the dataset.

The gain applied to frequency channel $f$ is
\begin{equation}
  B(f) = \max\Big( W(f)\,\big[\,1 + R(f) + S(f)\,\big],\; B_{\text{floor}} \Big),
  \label{eq:bandpass}
\end{equation}
where $W(f)$ is a Tukey (raised-cosine) taper that rolls the gain off over the
outer 8\% of the band at each edge, $R(f)$ is a slow ripple made of two to four
sine waves with periods between 0.4 and 1.5 times the band width and a total
amplitude of 15\%, and $S(f)$ is a standing wave with a period of 128 channels
and an amplitude of 5\%. $B_{\text{floor}} = 10^{-3}$ stops the gain from reaching
zero. Every parameter is randomly varied by up to $\pm25\%$ for each image and
all phases are random, so no two images share the same bandpass; otherwise a
model could simply memorise one curve.

An important detail is where the noise enters. In a real receiver chain the
sky signal and any RFI are amplified by the bandpass, but the electronic noise
of the amplifiers and digitiser is added afterwards and does not scale with
the gain. The spectrogram is therefore built as
\begin{equation}
  \text{spectrogram} = B(f)\,\big(\text{sky} + \text{RFI}\big) + n_{\text{rx}},
\end{equation}
where the receiver noise $n_{\text{rx}}$ has a standard deviation of 0.1 times
the median sky noise. If the noise were multiplied by the gain as well, it
would cancel out of the signal-to-noise ratio and the bandpass would make no
difference to how detectable the RFI is. With the noise added after the gain,
RFI near the band edges, where $B(f)$ is small, genuinely sinks into the noise.
The labelling rule of \cref{sec:synthetic_label} is applied after this step,
so such RFI is labelled clean. This removed 54{,}403 of the 37{,}073{,}066
injected RFI pixels (0.147\%), all near the band edges.

Each part of \cref{eq:bandpass} is standard: the Tukey window is a common
taper in signal processing\footcite{harris1978}, standing waves in radio
telescopes are well documented\footcite{popping2008}, and treating the instrument
as a gain applied to the incoming signal follows the radio-astronomy
measurement equation\footcite{hamaker1996}. The particular combination used here,
however, is my own construction and is not taken from a published simulator.

\section{HERA dataset}
\label{sec:hera}
The second dataset is simulated data for the Hydrogen Epoch of Reionization
Array (HERA), a radio interferometer in South Africa. It was produced by
Mesarcik et al.\footcite{mesarcik2022} and is publicly available on
Zenodo\footcite{mesarcik2022data}. The file contains 420 training and 140 test
spectrograms of 512$\times$512 pixels, each with a boolean RFI mask. Because
the data is simulated, the masks are exact. RFI covers 2.75\% of the pixels,
much less than the roughly 15\% in my synthetic dataset.

\subsection{Leakage between the training and test sets}
Before using the data I checked whether any test image also appears in the
training set, by comparing SHA-256 hashes of the raw pixel values. The
published split is not clean (\cref{tab:hera_leak}): both sets contain
internal duplicates, and 31 test images are exact copies of training
images---22\% of the test set. A model evaluated on this split is partly being
tested on images it has already seen. This is a property of the published
dataset, not something introduced by my processing.

\begin{table}[H]
  \centering
  \caption{Duplicate images in the published HERA split, found by comparing
    SHA-256 hashes of the raw pixels.}
  \label{tab:hera_leak}
  \begin{tabular}{lrr}
    \toprule
    Check & Published split & After deduplication \\
    \midrule
    Training images (unique / total) & 382 / 420 & 351 \\
    Test images (unique / total)     & 135 / 140 & 135 \\
    Test images also in training set & 31 (22\%) & 0 \\
    \bottomrule
  \end{tabular}
\end{table}

I removed the duplicates and the overlapping images, leaving 351 training and
135 test images with no overlap and no contradictory labels.

\subsection{Why HERA was not used further}
On the cleaned split my model (\cref{ch:hybrid}) reached a test \Fone{} of
0.9996 when trained from scratch, and 0.9964 when started from the weights
learned on the synthetic data. With both close to a perfect score, there is no
room for the dataset to show a difference between methods, so it cannot tell
us whether one approach is better than another. The leakage was not what made
the scores high: the score on the contaminated split, 0.9974, was almost the
same. I therefore treated HERA as too easy to be informative and moved to real
LOFAR data, which is described next. \tfunet{} was not run on HERA.

\section{LOFAR dataset}
\label{sec:lofar}
The third dataset is the most important one in this report, because it is the
only one made of real observations with labels checked by a person. It comes
from the Low-Frequency Array (LOFAR), a radio telescope network centred in the
Netherlands, and was released by Mesarcik et al.\footcite{mesarcik2022} on the same
Zenodo record as the HERA data\footcite{mesarcik2022data}.

\subsection{Origin and structure}
The data is a single 9.3\,GB Python pickle file. Loading it does not give
images but a list of four arrays:
\begin{itemize}
  \item training spectrograms, shape $(7500, 512, 512, 1)$;
  \item training masks for those spectrograms, same shape, true/false;
  \item test spectrograms, shape $(109, 512, 512, 1)$;
  \item test masks, same shape.
\end{itemize}
In each shape the first number is the number of spectrograms, the next two
are the rows and columns of one spectrogram, and the last is a single channel
(one value per pixel). According to the original paper, each spectrogram is a
512$\times$512 crop of a larger 599$\times$616 array.

The two sets of masks were made in different ways, and this matters for the
rest of the report. The training masks were produced automatically by
\aoflagger{}\footcite{offringa2012}, the standard flagging software for LOFAR.
The 109 test masks were labelled by hand by a human expert. The test set is
therefore the only real ground truth available.

\subsection{Problems found when opening the data}
\label{sec:lofar_audit}
Several properties of the file had to be found and handled before any model
could be trained on it.

\paragraph{Raw values.} The spectrograms are not normalised. Pixel values
range from 0 to $1.4\times10^{11}$, while the median is about $1.2\times10^{6}$;
the maximum is roughly 39{,}000 times the 99th percentile. When I first plotted
an image it appeared as a flat, single-coloured square, because a handful of
extremely bright pixels set the colour scale. Clipping and log-scaling the
values reveals the structure (\cref{fig:lofar_overview}).

\paragraph{Orientation.} In these arrays the rows are time and the columns are
frequency---the opposite of my synthetic data. I confirmed this from the masks:
persistent RFI stays at fixed frequencies, and 65\% of the flagged training
pixels (74\% in the test set) fall in just 16 of the 512 columns, against
about 25\% in the top 16 rows.

\paragraph{Test images inside the training set.} All 109 test spectrograms are
byte-for-byte identical to 109 of the training spectrograms. I verified this
by hashing every image and by direct element-wise comparison: the largest
pixel difference between each matched pair is exactly zero, while unrelated
images differ by about $10^{10}$. Only the masks differ---\aoflagger{} in the
training copy, the human expert in the test copy---and none of the 109 pairs
has identical masks. This is how the dataset was built, but it means that
training on all 7500 images would test the model on pictures it has already
seen. I removed these 109 images from the training set.

\paragraph{Dead images.} 35 training images are flagged as 100\% RFI; several
of these contain only zeros. They teach a model nothing useful and were also
removed. This leaves 7356 usable training images, of which 735 were held back
for validation and 6621 used for training.

\paragraph{Size.} The full file does not fit in memory alongside other work on
the 19\,GB machine used here, and loading it repeatedly crashed the
development environment. I converted it once into separate memory-mapped
arrays, which can be opened in 0.007\,s while using about 29\,MB of memory;
only the images actually needed are read from disk.

\begin{figure}[H]
  \centering
  \includegraphics[width=\textwidth]{fig_lofar_overview}
  \caption{Example LOFAR spectrograms (after clipping and log scaling for
    display) and their masks. Top two rows: training images with their
    \aoflagger{} masks; bottom two rows: test images with the human expert's
    masks. Vertical lines are RFI at fixed frequencies.}
  \label{fig:lofar_overview}
\end{figure}

\subsection{Why LOFAR is harder than the synthetic data}
\label{sec:lofar_harder}
LOFAR is harder than the synthetic data for three measurable reasons.

\paragraph{RFI is much rarer.} RFI covers 0.766\% of the pixels in the test
set (1.27\% in the training set), compared with about 15\% in the synthetic
data---roughly one RFI pixel for every 129 clean ones, against one in six. A
model that predicts ``no RFI'' everywhere would be 99.23\% accurate on the test
set while finding nothing, which is why accuracy is not used as a measure
in this report.

\paragraph{Much of the RFI is faint.} In the synthetic data every labelled
pixel is, by construction, above the noise (\cref{sec:synthetic_label}). In
LOFAR that is not true. To measure this I compared each RFI pixel with the
clean pixels of the same image: for each test image I found the 50th, 75th and
99th percentiles of its clean pixels, and placed each RFI pixel into a tier
according to where its brightness falls (\cref{tab:lofar_tiers}). Half of the
RFI is brighter than 99\% of the clean pixels and easy to see, but a fifth of
it is dimmer than a typical clean pixel. Those pixels carry no brightness
evidence at all; they can only be recognised from their surroundings.

\begin{table}[H]
  \centering
  \caption{Brightness of the human-labelled RFI pixels in the 109 LOFAR test
    images, measured against the clean pixels of the same image.}
  \label{tab:lofar_tiers}
  \begin{tabular}{llrr}
    \toprule
    Tier & Brightness relative to the image's clean pixels & RFI pixels & Share \\
    \midrule
    Bright    & above the 99th percentile          & 109{,}669 & 50.1\% \\
    Moderate  & between the 75th and 99th          &  41{,}263 & 18.9\% \\
    Faint     & between the 50th and 75th          &  24{,}113 & 11.0\% \\
    Invisible & below the median                   &  43{,}866 & 20.0\% \\
    \bottomrule
  \end{tabular}
\end{table}

\paragraph{The training labels are often wrong.} Because every test image also
appears in the training set, the \aoflagger{} mask and the human mask can be
compared on exactly the same pixels. Taking the human as correct, \aoflagger{}
has a precision of 0.560 and a recall of 0.580, an \Fone{} of 0.5698. Roughly
44\% of what any model learns from the training labels therefore disagrees
with the expert. As a check on my processing, Mesarcik et al.\ report the same
\aoflagger{} \Fone{} of 0.5698 in their paper\footcite{mesarcik2022}; my value,
0.5698403, agrees to four decimal places, which confirms that I am using the
same test data, labels and metric.
